%% clustering.tex -- shared sub-recipe. k-means and hierarchical linkage, the
%% two general-purpose grouping methods the library provides.

\begin{recipe}[label=sub:clustering]{Clustering Points and Trajectories}

\recipepurpose{Group things that behave alike: genes by expression vector,
  samples by whole time course, or anything else you can place in the space or
  give pairwise distances for. Two ordinary methods --- k-means and
  agglomerative linkage --- with the conventions Tensor Omics wraps them in.}

\nomaintainer

\begin{requiredinputs}
  \item For k-means: points as \code{(n\_dims, n\_points)}, one per column, and
        an initial centroid per cluster. There is no seeding routine ---
        supplying the centroids \emph{is} the seeding decision, and their count
        is $k$.
  \item For clustering whole trajectories: \code{(n\_factors, n\_samples,
        n\_timepoints)}, and one centroid per cluster of length
        \code{n\_factors}$\times$\code{n\_timepoints}.
  \item For linkage: a symmetric distance matrix with a zero diagonal and no
        negative entries. Computing it is not part of this step.
\end{requiredinputs}

\begin{prerequisites}
  \item None in the library. If you are clustering expression, normalize first
        --- otherwise the grouping follows overall magnitude.
\end{prerequisites}

\begin{workflow}
  \step{Decide what a \emph{point} is. Clustering genes, one column is a gene.
        Clustering samples by their evolution, one whole trajectory is one
        point --- that is
        \code{cluster\_factor\_trajectories\_k\_means}, and it is a different
        question from grouping the individual (sample, timepoint) states,
        which is plain \code{k\_means\_clustering} over the same values.}
  \step{For k-means, choose $k$ and the initial centroids together, since the
        array you pass carries both. Run \code{k\_means\_clustering}; the
        centroid array comes back holding the final centroids.}
  \step{For linkage, compute the distance matrix yourself and pass it with a
        criterion. The same matrix can be re-clustered under another criterion
        without recomputing it, and comes back unchanged either way.}
  \step{Read the merge sequence. Each step reports the two things merged, the
        height at which they met, and the size of the resulting cluster. Cut
        the tree at a height of your choosing --- the library does not cut it
        for you.}
\end{workflow}

\examplescript{python/examples/clustering.py}{%
  Clusters three planted blobs with k-means, groups four trajectories of which
  two rise and two fall, and prints an average-linkage tree merge by merge with
  the sign convention decoded. All data is synthetic and separable, so the
  expected grouping is obvious.}

\begin{codefragment}[language=Python,caption={Fragments to edit: k and its seeds, and the linkage criterion}]
import numpy as np
from tensor_omics import k_means_clustering, linkage_clustering

# --- k-means: the centroid array carries BOTH k and the seeds
points = np.asfortranarray(data)      # (n_dims, n_points)
centroids = np.asfortranarray(seeds)  # (n_dims, k)
res = k_means_clustering(points, centroids, max_iterations=300)
labels = res["labels"]                # 1-based, one per point
# centroids now holds the FINAL centroids: modified in place

# --- linkage: your distance matrix, their criterion --------
tree = linkage_clustering(
    np.asfortranarray(distances),     # symmetric, zero diagonal
    "average")                        # or "weighted" / "ward"

# positive entry = data point, NEGATIVE = the node made at
# that step. This is the opposite of R's hclust.
for k in range(len(tree["heights"])):
    i, j = tree["merge_i"][k], tree["merge_j"][k]
    height = tree["heights"][k]
\end{codefragment}

\begin{parameters}
  \param{centroids}{Both the number of clusters and where the search starts}{%
    $k$ is the number of columns. The values may be real data points or
    arbitrary --- but they must be distinct, which is \emph{not} checked}
  \param{max\_iterations}{When k-means gives up short of convergence}{Default
    \code{300}. Raise it only if you have reason to think it is stopping early;
    a k-means that needs many more usually has a seeding problem instead}
  \param{method}{The linkage criterion, i.e.\ how the distance between two
    clusters is defined}{\optn{average}, \optn{weighted} or \optn{ward}. Note
    what is \emph{not} here: \optn{single} and \optn{complete} are rejected}
\end{parameters}

\begin{expectedresults}
  On the example's three separated blobs, k-means recovers them exactly
  ($12/12/12$) and the final centroids sit within $0.1$ of the planted ones.
  The four trajectories split two and two, rising against falling. The linkage
  heights agree with a standard implementation --- they match
  \cmd{scipy.cluster.hierarchy.linkage} to floating-point tolerance for all
  three criteria --- and the last merge of the same eight points comes at
  $0.863$ under \optn{average}, $1.069$ under \optn{weighted} and $1.134$ under
  \optn{ward}, which is a useful reminder that heights are comparable within a
  criterion and not across them.
\end{expectedresults}

\begin{pitfall}[Pitfall: the seeding is yours, and duplicates are not caught]
  There is no k-means++ here, and no check that your initial centroids differ.
  Passing the same point $k$ times is accepted and returns $k$ populated
  clusters --- an arbitrary split of the data with no meaning at all. Seed from
  distinct data points, and since the outcome depends on that seeding, run it
  more than once from different seeds before believing a partition.
\end{pitfall}

\begin{pitfall}[Pitfall: the centroid array is overwritten]
  \code{centroids} is modified in place: the array you pass in comes back
  holding the \emph{final} centroids, which is how you retrieve them. Keep a
  copy first if you wanted your seeds afterwards, and do not reuse the same
  array across runs expecting the same starting point.
\end{pitfall}

\begin{pitfall}[Pitfall: the merge sign convention is the reverse of R's]
  In the merge arrays a \emph{positive} entry is a data point and a
  \emph{negative} entry is the inner node created at that step, so $-2$ means
  whatever step~2 produced. R's \cmd{hclust} uses exactly the opposite
  convention, negatives for singletons. Code ported from R will build a
  mirrored tree without any error being raised.
\end{pitfall}

\begin{pitfall}[Pitfall: expecting the distance matrix to be consumed]
  The matrix is declared as modified in place, and its lower triangle really is
  used as scratch --- but it is restored from the upper triangle before the
  call returns, on success and on error alike. There is no need to copy it
  defensively. It must, however, genuinely be a distance matrix: asymmetry, a
  negative entry or a non-zero diagonal each fail the call with
  \code{ToxInputError: invalid input arguments (argument 'distances')}.
\end{pitfall}

\begin{interpretation}
  The two methods answer differently shaped questions. k-means returns a
  partition at the $k$ you chose and says nothing about whether $k$ was right:
  every point gets a label, and a run on structureless data returns $k$ tidy
  clusters just the same. Linkage returns a whole hierarchy and leaves the cut
  to you, so the height at which groups merge is itself the evidence --- a
  merge far above the ones below it marks a real separation, and a tree that
  merges at steadily rising heights is telling you there is no natural number
  of groups.

  Whichever you use, the clustering describes the space you built, not biology
  directly. Cluster raw expression and you group by overall magnitude; cluster
  normalized vectors and you group by expression shape; cluster trajectories
  and you group by how things developed rather than where they ended. The
  choice made in the earlier recipes decides what your clusters mean.
\end{interpretation}

\begin{references}
  \reference{Lloyd (1982). Least squares quantization in PCM. \emph{IEEE
    Transactions on Information Theory} 28(2):129--137.}
  \reference{Ward (1963). Hierarchical grouping to optimize an objective
    function. \emph{Journal of the American Statistical Association}
    58(301):236--244.}
  \reference{M\"ullner (2011). Modern hierarchical, agglomerative clustering
    algorithms. \emph{arXiv}:1109.2378.}
\end{references}

\end{recipe}
